\textbf{Phần 1. Trắc nghiệm lựa chọn}
\setcounterpageref{ex}{0}

\begin{ex} %[5D4V20]
Tìm $x$, biết: $x + 14{,}65 = 32{,}4 - 8{,}25$.
\choice
{$8{,}5$}
{\True $9{,}5$}
{$10{,}5$}
{$9{,}25$}
\loigiai{
Ta có:
\begin{align*}
x + 14{,}65 &= 24{,}15\\
x &= 24{,}15 - 14{,}65\\
x &= 9{,}5
\end{align*}
\textbf{Đáp án của câu hỏi là B}
}
\end{ex}

\begin{ex} %[5D4V20]
Một cuộn dây điện dài $50\text{ m}$. Lần thứ nhất người ta cắt đi $16{,}75\text{ m}$, lần thứ hai cắt ít hơn lần thứ nhất $3{,}5\text{ m}$. Hỏi sau hai lần cắt, cuộn dây điện còn lại bao nhiêu mét?
\choice
{$19{,}5\text{ m}$}
{$20{,}5\text{ m}$}
{\True $20\text{ m}$}
{$29{,}75\text{ m}$}
\loigiai{
Lần thứ hai người ta cắt số mét dây điện là:
$16{,}75 - 3{,}5 = 13{,}25\text{ (m)}$.\\
Cả hai lần đã cắt tất cả số mét dây điện là:
$16{,}75 + 13{,}25 = 30\text{ (m)}$.\\
Sau hai lần cắt, cuộn dây còn lại số mét là:
$50 - 30 = 20\text{ (m)}$.\\
\textbf{Đáp án của câu hỏi là C}
}
\end{ex}

\begin{ex} %[5D4V20]
Tính nhanh giá trị của biểu thức: $45{,}68 - 18{,}35 - 11{,}65$.
\choice
{$16{,}68$}
{\True $15{,}68$}
{$15{,}35$}
{$16{,}35$}
\loigiai{
Ta nhóm hai số trừ lại:
$$45{,}68 - (18{,}35 + 11{,}65) = 45{,}68 - 30 = 15{,}68.$$
\textbf{Đáp án của câu hỏi là B}
}
\end{ex}

\begin{ex} %[5D4V20]
Hiệu của hai số bằng $18{,}5$. Nếu tăng số bị trừ thêm $4{,}2$ đơn vị và giảm số trừ đi $2{,}3$ đơn vị thì hiệu mới là bao nhiêu?
\choice
{$12\text{ đơn vị}$}
{$20{,}4\text{ đơn vị}$}
{\True $25\text{ đơn vị}$}
{$16{,}6\text{ đơn vị}$}
\loigiai{
Khi tăng số bị trừ thêm $4{,}2$ thì hiệu tăng thêm $4{,}2$.\\
Khi giảm số trừ đi $2{,}3$ thì hiệu tăng thêm $2{,}3$.\\
Hiệu mới bằng:
$$18{,}5 + 4{,}2 + 2{,}3 = 25.$$
\textbf{Đáp án của câu hỏi là C}
}
\end{ex}

\begin{ex} %[5D4V20]
Một thùng chứa dầu cân nặng $24{,}8\text{ kg}$. Sau khi người ta lấy ra một nửa lượng dầu trong thùng thì thùng và lượng dầu còn lại cân nặng $13{,}3\text{ kg}$. Hỏi riêng vỏ thùng cân nặng bao nhiêu ki-lô-gam?
\choice
{$1{,}5\text{ kg}$}
{\True $1{,}8\text{ kg}$}
{$2{,}3\text{ kg}$}
{$11{,}5\text{ kg}$}
\loigiai{
Nửa lượng dầu trong thùng nặng số ki-lô-gam là:
$24{,}8 - 13{,}3 = 11{,}5\text{ (kg)}$.\\
Toàn bộ lượng dầu trong thùng nặng số ki-lô-gam là:
$11{,}5 + 11{,}5 = 23\text{ (kg)}$.\\
Vỏ thùng rỗng nặng số ki-lô-gam là:
$24{,}8 - 23 = 1{,}8\text{ (kg)}$.\\
\textbf{Đáp án của câu hỏi là B}
}
\end{ex}

\begin{ex} %[5D4V20]
Khi trừ một số tự nhiên cho một số thập phân có một chữ số ở phần thập phân, một bạn học sinh đã quên viết dấu phẩy của số trừ nên thực hiện phép trừ như trừ hai số tự nhiên và được kết quả là $145$. Hiệu đúng phải là $381{,}7$. Số thập phân cần tìm là
\choice
{$23{,}6$}
{\True $26{,}3$}
{$27{,}4$}
{$25{,}2$}
\loigiai{
Khi bỏ quên dấu phẩy ở số thập phân có một chữ số thập phân thì số trừ đã gấp lên $10$ lần.\\
Do đó, hiệu giảm đi một lượng bằng $9$ lần số thập phân đó.\\
Chín lần số thập phân đó là:
$381{,}7 - 145 = 236{,}7$.\\
Số thập phân đó là:
$236{,}7 : 9 = 26{,}3$.\\
\textbf{Đáp án của câu hỏi là B}
}
\end{ex}

\begin{ex} %[5D4V20]
An và Bình cân nặng $65{,}4\text{ kg}$; Bình và Cúc cân nặng $68{,}2\text{ kg}$. Hỏi Cúc nặng hơn An bao nhiêu ki-lô-gam?
\choice
{$3{,}2\text{ kg}$}
{$2{,}6\text{ kg}$}
{\True $2{,}8\text{ kg}$}
{$1{,}8\text{ kg}$}
\loigiai{
Cúc nặng hơn An số ki-lô-gam là:
$(\text{Bình} + \text{Cúc}) - (\text{An} + \text{Bình}) = 68{,}2 - 65{,}4 = 2{,}8\text{ (kg)}$.\\
\textbf{Đáp án của câu hỏi là C}
}
\end{ex}

\begin{ex} %[5D4V20]
Giá trị của biểu thức $100 - (38{,}6 + 21{,}4) - 15{,}75$ là
\choice
{$25{,}25$}
{\True $24{,}25$}
{$23{,}75$}
{$24{,}75$}
\loigiai{
Ta có:
$$100 - (38{,}6 + 21{,}4) - 15{,}75 = 100 - 60 - 15{,}75 = 40 - 15{,}75 = 24{,}25.$$
\textbf{Đáp án của câu hỏi là B}
}
\end{ex}

\begin{ex} %[5D4V20]
Một thửa ruộng hình chữ nhật có chu vi $64\text{ m}$, chiều dài hơn chiều rộng $8{,}6\text{ m}$. Chiều rộng của thửa ruộng đó là
\choice
{$12{,}7\text{ m}$}
{\True $11{,}7\text{ m}$}
{$20{,}3\text{ m}$}
{$10{,}7\text{ m}$}
\loigiai{
Nửa chu vi của thửa ruộng là:
$64 : 2 = 32\text{ (m)}$.\\
Chiều rộng của thửa ruộng là:
$(32 - 8{,}6) : 2 = 23{,}4 : 2 = 11{,}7\text{ (m)}$.\\
\textbf{Đáp án của câu hỏi là B}
}
\end{ex}

\begin{ex} %[5D4V20]
Cho các chữ số $2, 5, 8$. Lập số thập phân lớn nhất và số thập phân nhỏ nhất có một chữ số ở phần nguyên và hai chữ số ở phần thập phân từ ba chữ số trên (mỗi chữ số dùng một lần). Hiệu của hai số lập được là
\choice
{$5{,}84$}
{$6{,}04$}
{\True $5{,}94$}
{$5{,}74$}
\loigiai{
Số lớn nhất lập được là: $8{,}52$.\\
Số nhỏ nhất lập được là: $2{,}58$.\\
Hiệu của hai số là:
$8{,}52 - 2{,}58 = 5{,}94$.\\
\textbf{Đáp án của câu hỏi là C}
}
\end{ex}

\begin{ex} %[5D4V20]
Mẹ mang $100\,000$ đồng đi chợ. Mẹ mua thịt hết $52\,500$ đồng và mua rau củ quả hết $27\,500$ đồng. Hỏi sau khi mua xong, mẹ còn lại bao nhiêu tiền?
\choice
{$30\,000\text{ đồng}$}
{\True $20\,000\text{ đồng}$}
{$25\,000\text{ đồng}$}
{$15\,000\text{ đồng}$}
\loigiai{
Tổng số tiền mẹ đã chi tiêu là:
$52\,500 + 27\,500 = 80\,000\text{ (đồng)}$.\\
Số tiền mẹ còn lại là:
$100\,000 - 80\,000 = 20\,000\text{ (đồng)}$.\\
\textbf{Đáp án của câu hỏi là B}
}
\end{ex}

\begin{ex} %[5D4V20]
Có ba đoạn dây. Đoạn thứ nhất dài $14{,}5\text{ m}$. Đoạn thứ hai ngắn hơn đoạn thứ nhất $2{,}8\text{ m}$. Đoạn thứ ba ngắn hơn tổng chiều dài hai đoạn đầu là $5{,}4\text{ m}$. Chiều dài đoạn thứ ba là
\choice
{$21{,}8\text{ m}$}
{\True $20{,}8\text{ m}$}
{$19{,}8\text{ m}$}
{$26{,}2\text{ m}$}
\loigiai{
Đoạn thứ hai dài số mét là:
$14{,}5 - 2{,}8 = 11{,}7\text{ (m)}$.\\
Tổng chiều dài hai đoạn đầu là:
$14{,}5 + 11{,}7 = 26{,}2\text{ (m)}$.\\
Đoạn thứ ba dài số mét là:
$26{,}2 - 5{,}4 = 20{,}8\text{ (m)}$.\\
\textbf{Đáp án của câu hỏi là B}
}
\end{ex}

\textbf{Phần 2. Trắc nghiệm đúng sai}
\setcounterpageref{ex}{0}

\begin{ex} %[5D4V20]
Cho biểu thức $A = 78{,}5 - 24{,}65 - 15{,}35$ và $B = 60 - (12{,}8 + 18{,}7)$.
\choiceTF
{\True Giá trị của biểu thức $A$ bằng $38{,}5$}
{\True Giá trị của biểu thức $B$ bằng $28{,}5$}
{\True Hiệu $A - B = 10$}
{Tổng $A + B = 66$}
\loigiai{
\begin{itemchoice}
\item Đúng. Áp dụng quy tắc trừ một tổng: $A = 78{,}5 - (24{,}65 + 15{,}35) = 78{,}5 - 40 = 38{,}5$.
\item Đúng. Ta có $B = 60 - 31{,}5 = 28{,}5$.
\item Đúng. Hiệu $A - B = 38{,}5 - 28{,}5 = 10$.
\item Sai. Tổng $A + B = 38{,}5 + 28{,}5 = 67 
eq 66$.
\end{itemchoice}
}
\end{ex}

\begin{ex} %[5D4V20]
Một cửa hàng gạo có $1{,}5$ tấn gạo. Ngày thứ nhất cửa hàng bán được $380\text{ kg}$ gạo. Ngày thứ hai bán được nhiều hơn ngày thứ nhất $90\text{ kg}$ gạo.
\choiceTF
{\True Đổi $1{,}5\text{ tấn} = 1\,500\text{ kg}$}
{\True Ngày thứ hai cửa hàng bán được $470\text{ kg}$ gạo}
{Sau hai ngày cửa hàng đã bán được tất cả $800\text{ kg}$ gạo}
{\True Sau hai ngày cửa hàng còn lại $650\text{ kg}$ gạo}
\loigiai{
\begin{itemchoice}
\item Đúng. Vì $1\text{ tấn} = 1\,000\text{ kg}$ nên $1{,}5\text{ tấn} = 1\,500\text{ kg}$.
\item Đúng. Ngày thứ hai bán được: $380 + 90 = 470\text{ (kg)}$.
\item Sai. Cả hai ngày bán được: $380 + 470 = 850\text{ (kg)}$.
\item Đúng. Số gạo còn lại là: $1\,500 - 850 = 650\text{ (kg)}$.
\end{itemchoice}
}
\end{ex}

\textbf{Phần 3. Trắc nghiệm trả lời ngắn}
\setcounterpageref{ex}{0}

\begin{ex} %[5D4V20]
Tìm số lượng các số tự nhiên $x$ thoả mãn điều kiện: $12{,}5 - x > 7{,}8$.
\shortans{5}
\loigiai{
Ta có: $12{,}5 - x > 7{,}8 \implies x < 12{,}5 - 7{,}8 = 4{,}7$.\\
Vì $x$ là số tự nhiên nên $x$ nhận các giá trị: $0; 1; 2; 3; 4$.\\
Vậy có tất cả 5 giá trị tự nhiên thoả mãn.
}
\end{ex}

\begin{ex} %[5D4V20]
Hiệu của hai số là $16{,}5$. Nếu tăng số bị trừ thêm $5{,}5$ đơn vị và giảm số trừ đi $4$ đơn vị thì hiệu mới bằng bao nhiêu?
\shortans{26}
\loigiai{
Khi số bị trừ tăng thêm $5{,}5$ và số trừ giảm đi $4$ thì hiệu mới tăng thêm:
$$5{,}5 + 4 = 9{,}5.$$
Hiệu mới là:
$$16{,}5 + 9{,}5 = 26.$$
}
\end{ex}

\begin{ex} %[5D4V20]
Một thùng chứa $20\text{ l}$ nước hoa quả. Người ta rót ra 6 chai, mỗi chai $1{,}25\text{ l}$, và rót vào một bình lớn $7{,}5\text{ l}$. Hỏi trong thùng còn lại bao nhiêu lít nước hoa quả?
\shortans{5}
\loigiai{
Số lít nước hoa quả rót vào 6 chai là:
$1{,}25 \times 6 = 7{,}5\text{ (l)}$.\\
Tổng số lít nước hoa quả đã lấy ra là:
$7{,}5 + 7{,}5 = 15\text{ (l)}$.\\
Số lít nước hoa quả còn lại trong thùng là:
$20 - 15 = 5\text{ (l)}$.
}
\end{ex}

\begin{ex} %[5D4V20]
Một cuộn vải dài $45\text{ m}$. Lần thứ nhất người ta cắt bán $14{,}8\text{ m}$. Lần thứ hai cắt bán ít hơn lần thứ nhất $2{,}6\text{ m}$. Hỏi sau hai lần bán, cuộn vải còn lại bao nhiêu mét?
\shortans{18}
\loigiai{
Lần thứ hai cắt bán số mét vải là:
$14{,}8 - 2{,}6 = 12{,}2\text{ (m)}$.\\
Cả hai lần đã bán số mét vải là:
$14{,}8 + 12{,}2 = 27\text{ (m)}$.\\
Cuộn vải còn lại số mét là:
$45 - 27 = 18\text{ (m)}$.
}
\end{ex}

\begin{ex} %[5D4V20]
Một ô tô chở hàng có khối lượng cả xe và hàng là $5{,}2$ tấn. Sau khi dỡ xuống $\dfrac{3}{5}$ số hàng thì khối lượng cả xe và lượng hàng còn lại là $3{,}4$ tấn. Hỏi riêng xe ô tô (khi không chở hàng) có khối lượng là bao nhiêu tạ?
\shortans{22}
\loigiai{
Khối lượng của $\dfrac{3}{5}$ số hàng được dỡ xuống là:
$5{,}2 - 3{,}4 = 1{,}8\text{ (tấn)}$.\\
Tổng khối lượng của tất cả hàng hoá ban đầu là:
$1{,}8 : 3 \times 5 = 3\text{ (tấn)}$.\\
Khối lượng của xe khi không chở hàng là:
$5{,}2 - 3 = 2{,}2\text{ (tấn)}$.\\
Đổi $2{,}2\text{ tấn} = 22\text{ tạ}$.
}
\end{ex}

\begin{ex} %[5D4V20]
Ba số có tổng bằng $100$. Số thứ nhất là $34{,}6$. Số thứ hai hơn số thứ ba là $8{,}6$. Tìm số thứ hai.
\shortans{37}
\loigiai{
Tổng của số thứ hai và số thứ ba là:
$100 - 34{,}6 = 65{,}4$.\\
Số thứ hai là:
$(65{,}4 + 8{,}6) : 2 = 74 : 2 = 37$.
}
\end{ex}
